\section{回到拓扑哈密顿动力学}

第 6 章引入了拓扑哈密顿流及其相伴哈密顿量的概念。本节说明如何把各种谱量延拓到
\emph{拓扑}哈密顿道路。

\subsection{拓扑哈密顿量的谱不变量}

本小节把\emph{哈密顿道路}的谱不变量定义及基本性质延拓到拓扑哈密顿范畴。

为完成这一延拓，关键是像 (Oh05c) 那样在哈密顿道路空间
\[
 \mathcal P^{\mathrm{ham}}(\Symp(M,\omega),id)
\]
的层面，而不只在
$\Ham(M,\omega)$ 的覆盖空间上给出定义。此外，第 6.3 节的拓扑哈密顿量唯一性定理
对延拓到这一 $C^0$ 范畴至关重要。

先用哈密顿道路 $\phi_H$ 改写谱数 $\rho(H;a)$：
\[
 \rho(\lambda;a):=\rho(\underline H;a)
 =\inf_{\alpha\in\ker\partial_H}\{\lambda_{\underline H}(\alpha)\mid[\alpha]=a^\flat\}.
\]
对光滑哈密顿量 $H$，这只是换一种写法。为了用一列 $H_i$（其哈密顿流
$\phi_{H_i}$ 一致收敛到给定的拓扑哈密顿道路 $\lambda$）所对应的
$\rho(\underline H_i;a)$ 之极限来延拓定义，需要证明
$\lim_{i\to\infty}\rho(\underline H_i;a)$ 与这类序列 $H_i$ 的选择无关。

\begin{theorem}
\label{thm:22-topological-spectral-invariant}
对带归一化哈密顿量 $H$ 的光滑哈密顿道路 $\phi_H$，定义
\[
 \rho(\phi_H;a)=\rho(\underline H;a).
\]
则映射 $\rho_a:\phi_H\mapsto\rho(\phi_H;a)$ 延拓为（在哈密顿拓扑下）连续函数
\[
 \rho_a=\rho(\,\cdot\,;a):
 \mathcal P^{\mathrm{ham}}(\operatorname{Sympeo}(M,\omega),id)
 \longrightarrow\mathbb R,
\]
并满足三角不等式
\begin{equation}
 \rho(\lambda\mu;a\cdot b)\leq\rho(\lambda;a)+\rho(\mu;b).
 \label{eq:22-topological-triangle}
\end{equation}
\end{theorem}

\begin{proof}
第一条陈述是
\[
 \rho_a:\mathcal P^{\mathrm{ham}}(\Symp(M,\omega),id)\longrightarrow\mathbb R
\]
的哈密顿连续性与定理 6.3.3 的直接推论：对任意拓扑哈密顿道路 $\lambda$，定义
\begin{equation}
 \rho(\lambda;a)=\lim_{i\to\infty}\rho(\underline H_i;a),
 \label{eq:22-topological-spectral-limit}
\end{equation}
其中 $H_i$ 是任意满足 $\phi_{H_i}\to\lambda$ 且在 $L^{(1,\infty)}$ 中收敛的序列。
唯一性定理保证该定义良定。此外，命题 21.3.6 蕴含此延拓在
$\mathcal P^{\mathrm{ham}}(\operatorname{Sympeo}(M,\omega),id)$ 上连续。对光滑情形的
三角不等式取极限，立刻得到~\eqref{eq:22-topological-triangle}。
\end{proof}

现在专注于 $1\in QH^*(M)$ 的不变量 $\rho(\lambda;1)$。回忆，对光滑哈密顿道路
$\lambda=\phi_H$，此前定义了
\[
 \gamma(H)=\rho(H;1)+\rho(\overline H;1).
\]
为免混淆，当把 $\gamma$ 看作
$\mathcal P^{\mathrm{ham}}(\Symp(M,\omega),id)$ 上的函数时，记作
$\operatorname{norm}_\gamma$。

\begin{definition}[谱伪范数]
\label{def:22-spectral-pseudonorm}
令 $\lambda\in\mathcal P^{\mathrm{ham}}(\Symp(M,\omega),id)$，并令 $H$ 为满足
$\lambda=\phi_H$ 的哈密顿量。定义函数
\[
 \operatorname{norm}_\gamma:
 \mathcal P^{\mathrm{ham}}(\Symp(M,\omega),id)\longrightarrow\mathbb R_+
\]
为 $\operatorname{norm}_\gamma(\lambda)=\gamma(H)$。称
$\operatorname{norm}_\gamma(\lambda)$ 为 $\lambda$ 的\emph{谱伪范数}。
\end{definition}

此前已经证明它非负，并且只依赖于 $\lambda=\phi_H$ 在 $\Ham(M,\omega)$ 中的道路同伦类。
拓扑哈密顿量的唯一性再次使我们能把定义延拓到拓扑哈密顿道路。

\begin{proposition}
\label{prop:22-topological-spectral-pseudonorm}
谱伪范数 $\gamma$ 延拓为连续函数
\[
 \operatorname{norm}_{\overline\gamma}:
 \mathcal P^{\mathrm{ham}}(\operatorname{Sympeo}(M,\omega),id)
 \longrightarrow\mathbb R,
\]
其定义为
\begin{equation}
 \operatorname{norm}_{\overline\gamma}(\lambda)
 =\lim_{i\to\infty}\operatorname{norm}_\gamma(\phi_{H_i}),
 \label{eq:22-topological-spectral-pseudonorm}
\end{equation}
其中 $\phi_{H_i}\to\lambda$ 是 $L^{(1,\infty)}$-哈密顿拓扑中的某个（因而任意）序列。
\end{proposition}

回忆，对光滑哈密顿量 $H$，每个
$\rho(\phi_H;a)=\rho(H;a)$ 都与哈密顿方程
$\dot x=X_H(t,x)$ 的一个周期轨道相伴，并对应于该周期轨道的作用量；至少这对有理辛
流形上的\emph{任意}哈密顿量成立（证明见 (Oh05c, Oh09a)），也对一般无理辛流形上的\emph{非退化}
哈密顿量成立（见 (Oh09a), (Ush08)）。由此看来，下述问题似乎至关重要。

\begin{question}
\label{ques:22-topological-pseudonorm-dynamics}
\stfhIndex{I178}{Puweifanshu}{谱伪范数（spectral pseudonorm）}
延拓后的谱伪范数 $\operatorname{norm}_{\overline\gamma}(\lambda)$ 对拓扑哈密顿流的动力学
有何含义？
\end{question}

下面是 (Oh05d) 中谱位移能定义在拓扑哈密顿范畴内的对应版本。

\begin{definition}[谱位移能]
\label{def:22-topological-spectral-displacement-energy}
令 $A\subset M$ 为紧子集。定义 $A$ 的谱位移能 $e_{\overline\gamma}(A)$ 为
\[
 e_{\overline\gamma}(A)
 =\inf_\lambda\{\operatorname{norm}_{\overline\gamma}(\lambda)
 \mid A\cap\lambda(1)(A)=\varnothing,
 \ \lambda\in\mathcal P_\infty^{\mathrm{ham}}
 (\operatorname{Sympeo}(M,\omega),id)\}.
\]
\end{definition}

展开哈密顿同胚与谱位移能的定义，还得到下述定理。

\begin{theorem}
\label{thm:22-spectral-displacement-agreement}
对任意紧子集 $A\subset M$，有
$e_{\overline\gamma}(A)=e_\gamma(A)$。
\end{theorem}

基于该定理，即使在拓扑哈密顿范畴中，也把 $A$ 的谱位移能简记为 $e_\gamma(A)$。于是
有下述定理。

\begin{theorem}
\label{thm:22-spectral-displacement-invariance}
对每个 $\psi\in\operatorname{Sympeo}(M,\omega)$，
\[
 e_\gamma(A)=e_\gamma(\psi(A)).
\]
\end{theorem}

\begin{proof}
注意 $h(A)\cap A=\varnothing$ 当且仅当
$\psi h\psi^{-1}(\psi(A))\cap\psi(A)=\varnothing$。此外，
$h\in\Hameo(M,\omega)$ 当且仅当
$\psi h\psi^{-1}\in\Hameo(M,\omega)$。
\end{proof}

接着考虑谱范数
\[
 \gamma(\phi):=\inf_{\operatorname{ev}_1(\lambda)=\phi}
 \operatorname{norm}_\gamma(\lambda)
 =\inf_{H\mapsto\phi}\{\rho(H;1)+\rho(\overline H;1)\}.
\]
下述定义把谱范数延拓到 $\Hameo(M,\omega)$。

\begin{definition}[谱范数]
\label{def:22-generalized-spectral-norm}
令 $h\in\Hameo(M,\omega)$，考虑满足
$\overline{\operatorname{ev}}_1(\lambda)=h$ 的拓扑哈密顿道路
$\lambda\in\mathcal P_\infty^{\mathrm{ham}}
(\operatorname{Sympeo}(M,\omega),id)$。若
$\overline{\operatorname{ev}}_1(\lambda)=h$，简记为 $\lambda\mapsto h$。定义
\footnote{译者注：原书定义首句把端点评价条件印作
$\overline{\operatorname{ev}}_1(0)=h$，但下一句及式 (22.6.70) 都使用
$\overline{\operatorname{ev}}_1(\lambda)=h$；译文按后两处改正。}
\begin{equation}
 \overline\gamma(h)=\inf_\lambda
 \{\operatorname{norm}_{\overline\gamma}(\lambda)
 \mid \lambda\in\mathcal P_\infty^{\mathrm{ham}}
 (\operatorname{Sympeo}(M,\omega),id),
 \ \overline{\operatorname{ev}}_1(\lambda)=h\}.
 \label{eq:22-generalized-spectral-norm}
\end{equation}
\end{definition}

下述结果在拓扑哈密顿语境中建立了不变范数全部性质的对应版本。

\begin{theorem}
\label{thm:22-generalized-spectral-norm-properties}
广义谱函数
$\overline\gamma:\Hameo(M,\omega)\to\mathbb R_+$ 满足不变范数的全部性质。
\end{theorem}

一旦证明光滑情形的下述连续性引理，证明便与霍费尔范数的证明本质相同。

\begin{lemma}
\label{lem:22-spectral-norm-hofer-continuity}
函数 $\gamma:\Ham(M,\omega)\to\mathbb R_+$ 在霍费尔拓扑中连续，因而也在
$\Ham(M,\omega)$ 的哈密顿拓扑中连续。
\end{lemma}

\begin{proof}
设 $H\mapsto\phi$、$K\mapsto\psi$。$\gamma$ 的三角不等式及
$\gamma(\phi)\leq\lVert\phi\rVert$ 蕴含
\[
 |\gamma(\phi)-\gamma(\psi)|
 \leq\gamma(\phi^{-1}\psi)
 \leq\lVert\phi^{-1}\psi\rVert
 =d_{\mathrm{Hofer}}(\phi,\psi).
\]
\end{proof}

由此产生如下基础问题。

\begin{question}
\label{ques:22-spectral-norm-continuity}
\begin{enumerate}[label=(\arabic*)]
\item $\gamma$（或 $\overline\gamma$）在 $C^0$-拓扑中连续吗？
\item 恒等式
\begin{equation}
 \overline\gamma|_{\Ham(M,\omega)}=\gamma
 \label{eq:22-spectral-norm-restriction}
\end{equation}
成立吗？
\end{enumerate}
\end{question}

最近，塞法迪尼 (Sey13a) 综合 $C^0$ 辛拓扑与哈密顿动力学中的多种要素，对二维曲面
肯定回答了第一个问题。它又在二维保面积动力系统中导出一个有趣推论 (Sey13b)。

\subsection{\texorpdfstring{$\mathcal P^{\mathrm{ham}}(\operatorname{Sympeo}(S^2),id)$}
{P-ham(Sympeo(S2),id)} 上的卡拉比拟同态}

本节余下部分把讨论限制在带标准辛形式 $\omega_{S^2}$ 的球面 $S^2$。在记号中省略
$\omega_{S^2}$，以
\[
 \mathcal P^{\mathrm{ham}}(\operatorname{Sympeo}(S^2),id)
 \quad\text{和}\quad \Hameo(S^2)
\]
分别简记 $S^2$ 上的拓扑哈密顿道路群与哈密顿同胚群，其他记号类似。

由上述定义及 $\rho(\,\cdot\,;1)$ 的哈密顿连续性，
$\mu^{\mathrm{path}}$ 也哈密顿连续。$\mu^{\mathrm{path}}$ 的卡拉比性质立即蕴含：
把 $\mu$ 限制到
$\mathcal P^{\mathrm{ham}}(\Symp_U(S^2),id)$ 后具有如下同态性质。

\begin{corollary}
\label{cor:22-local-path-calabi-homomorphism}
设 $U$ 是 $S^2$ 的开子集，$\overline U$ 在 $S^2$ 上可移，并令
\[
 \lambda_1,\lambda_2\in\mathcal P^{\mathrm{ham}}(\Symp_U(S^2),id).
\]
则
\[
 \mu^{\mathrm{path}}(\lambda_1\lambda_2)
 =\mu^{\mathrm{path}}(\lambda_1)+\mu^{\mathrm{path}}(\lambda_2).
\]
\end{corollary}

上述全部讨论都可延拓到拓扑哈密顿道路层面。只要知道下列事实，推广便立即得到：
\begin{enumerate}[label=(\arabic*)]
\item 对任意闭辛流形，$\rho(\,\cdot\,;1)$ 已延拓到
$\mathcal P^{\mathrm{ham}}(\operatorname{Sympeo}(M,\omega),id)$；特别地，第 12.6 节
处理了 $(S^2,\omega_{S^2})$\footnote{译者注：原书此处交叉引用为第 12.6 节；但所指的正是本节前述拓扑延拓，
且本书第 12.6 节并不讨论该问题，疑应为第 22.6 节。译文保留原编号并提示。}；
\item 此外，该延拓哈密顿连续，即在哈密顿拓扑中连续。
\end{enumerate}
综上得到下述定理。

\begin{theorem}
\label{thm:22-topological-calabi-quasimorphism}
卡拉比拟同态
\[
 \mu^{\mathrm{path}}:
 \mathcal P^{\mathrm{ham}}(\Symp(S^2),id)\longrightarrow\mathbb R
\]
可延拓为拟同态
\[
 \overline\mu^{\mathrm{path}}:
 \mathcal P^{\mathrm{ham}}(\operatorname{Sympeo}(S^2),id)
 \longrightarrow\mathbb R,
\]
并满足命题~\ref{prop:22-descend-calabi-quasimorphism} 与卡拉比性质的所有对应版本。
\end{theorem}

下面提出一个猜想；我们相信，它在研究 $S^2$ 保面积群的单纯性问题时将起重要作用。
回忆 (EnP03) 已由恩托夫和波尔捷罗维奇对 $S^2$ 上\emph{光滑}哈密顿微分同胚群
$\Ham(S^2)$ 证明了相应事实。

\begin{conjecture}
\label{conj:22-descend-topological-calabi}
令
$\overline\mu^{\mathrm{path}}:
\mathcal P^{\mathrm{ham}}(\operatorname{Sympeo}(S^2),id)\to\mathbb R$
为~\eqref{eq:22-path-homogeneous-quasimorphism} 给出的齐次卡拉比拟同态之上述延拓。
它下降为齐次拟同态
$\overline\mu:\Hameo(S^2)\to\mathbb R$，满足
\begin{equation}
 \overline\mu^{\mathrm{path}}
 =\overline\mu\circ\overline{\operatorname{ev}}_1.
 \label{eq:22-topological-calabi-descent}
\end{equation}
特别地，只要
$\lambda\in\mathcal P^{\mathrm{ham}}(\operatorname{Sympeo}(S^2),id)$，
$\overline\mu^{\mathrm{path}}(\lambda)$ 便只依赖于 $\lambda$ 的时一映射 $\lambda(1)$。
此外，对任意 $\phi\in\Hameo(S^2)$，只要存在支集在 $\operatorname{Int}(D^+)$ 中、
且满足 $\lambda(1)=\phi$ 的拓扑哈密顿道路 $\lambda$，就有
\begin{equation}
 \overline\mu(\phi)=\overline{\operatorname{Cal}}_{D^+}(\phi).
 \label{eq:22-topological-calabi-property}
\end{equation}
\end{conjecture}

猜想~\ref{conj:22-descend-topological-calabi} 与
$\overline\mu^{\mathrm{path}}$ 的卡拉比性质会立即推出下述猜想。

\begin{conjecture}
\label{conj:22-calabi-homeomorphism-extension}
卡拉比同态
$\operatorname{Cal}:\Ham(D^2,\partial D^2)\to\mathbb R$ 可延拓为在哈密顿拓扑中连续的
同态
\[
 \overline{\operatorname{Cal}}:
 \Hameo(D^2,\partial D^2)\longrightarrow\mathbb R.
\]
\end{conjecture}
